We relate the unitary spherical pair to Shalika periods with
coefficient representations through the similitude theta correspondence. For ordinary
unitary groups, the theta-period identity is given in
Gan~\cite[\S2]{GanPeriods} and Chen--Gan
\cite[\S3, Proposition~3.4 and Theorem~3.5]{ChenGan}.
Here the multiplier coordinate enters the rational orbit and the
stabilizer action. Theorem~\ref{thm:main-unitary} then bounds the
resulting Hom space.
\label{gu:sec:setup}

\subsection{The unitary similitude dual pair}
We use the notation of the introduction and regard vectors as column
vectors, with groups acting on the left. Let $J_4$ be the antidiagonal
identity matrix and set
\[
 V=E^4,\qquad h_V(v,v')=\bar v'^{\,t}J_4v.
\]
Let $W=X\oplus Y$, with $X=Y=E^2$, and put
\begin{equation}\label{gu:eq:W-form}
 h_W((x,y),(x',y'))=x^t\bar y'-y^t\bar x',
 \qquad J_W=\begin{pmatrix}0&I_2\\-I_2&0\end{pmatrix}.
\end{equation}
Thus
\[
 G_\theta=\GU(V)=\{g:g^\dagger J_4g=\lambda_{G_\theta}(g)J_4\},\qquad
 G'_\theta=\GU(W)=\{h:h^\dagger J_Wh=\lambda_{G'_\theta}(h)J_W\}.
\]
Write $G_\theta^1=U(V)$ and $(G'_\theta)^1=U(W)$ for their kernels. When the group is clear, we abbreviate its multiplier to $\lambda$. Multiplication of $h_W$ by a nonzero trace-zero element of $E$ identifies $G'_\theta$ with the split Hermitian similitude group $\GU_{2,2}(F)$, preserving its multiplier. The skew-Hermitian realization is adapted to the dual pair.

Fix a hyperbolic Hermitian form on $X$,
\[
 B(x,x')=\bar x'^{\,t}\mathsf Bx,\qquad \mathsf B^\dagger=\mathsf B,
\]
and put $K=\GU(B)$. One may take $\mathsf B=J_2$. Keeping $\mathsf B$ in the formulas allows arbitrary hyperbolic bases.

The conjugate space and its form are
\[
 X^c=X\otimes_{E,\bar{\ }}E,
 \qquad x\otimes z=\bar z x\otimes1,
 \qquad B^c(x\otimes1,x'\otimes1)=B(x',x).
\]
For $a\in\GL(X)$ define $\iota(a)(x\otimes1)=ax\otimes1$. In conjugate bases, $\iota(a)$ has matrix $\bar a$, and $B^c$ has matrix $\overline{\mathsf B}$. Consequently
\begin{equation}\label{gu:eq:iota-det}
 \iota:K\xrightarrow{\sim}\GU(B^c),\qquad
 \lambda(\iota(a))=\lambda_B(a),\qquad
 \det_{X^c}\iota(a)=\overline{\det_X a}.
\end{equation}
These are the conjugate-space conventions of~\cite[Sections~2.2--2.4]{ChenGan}.

\subsection{The Shalika character and subgroup}
The Siegel Levi and unipotent radical in $G'_\theta$ are
\begin{equation}\label{gu:eq:siegel-coordinates}
 m(a,t)=\begin{pmatrix}a&0\\0&t\bar a^{-t}\end{pmatrix},
 \qquad
 n(u)=\begin{pmatrix}I_2&u\\0&I_2\end{pmatrix},
 \quad u^\dagger=u.
\end{equation}
Here $a\in\GL_2(E)$, $t\in F^\times$, and $\lambda(m(a,t))=t$. Let $P^{\GU}$ be the Siegel parabolic with this Levi and unipotent radical $N(X)$. The form $B$ corresponds to the opposite-unipotent map $b_B:X\to Y$ with matrix
\[
 b_B=-\mathsf B.
\]
Indeed, $h_W(b_Bx,x')=-\bar x'^{\,t}b_Bx=B(x,x')$.

We use the additive character $\psi$ fixed in the introduction and
write $\tr_E$ for matrix trace over $E$. Our normalization is
\begin{equation}\label{gu:eq:psiB}
 \psi_B(n(u))
 =\psi\left(-\tfrac12\tr_E(u\mathsf B)\right)
 =\psi\left(-\tfrac14\Tr_{E/F}\tr_E(u\mathsf B)\right).
\end{equation}
Here $\tr_E(u\mathsf B)\in F$, because both matrices are Hermitian.

Write $\operatorname{Herm}_2(E)$ for the space of Hermitian $2$ by $2$ matrices. The trace pairing $(u,D)\mapsto\tr_E(uD)$ on this space
is nondegenerate: diagonal test matrices recover the diagonal entries,
and the off-diagonal entries are detected by the field-trace pairing.
Since
\begin{equation}\label{gu:jac:levi-conjugation}
 m(a,t)n(u)m(a,t)^{-1}=n(t^{-1}aua^\dagger),
\end{equation}
we have
\[
 \psi_B(m(a,t)n(u)m(a,t)^{-1})
 =\psi\left(-\tfrac12\tr_E(u\,t^{-1}a^\dagger\mathsf Ba)\right).
\]
Nondegeneracy and the nontriviality of $\psi$ show that the character
stabilizer is given by
\begin{equation}\label{gu:jac:character-stabilizer}
 a^\dagger\mathsf Ba=t\mathsf B.
\end{equation}
Thus
\begin{equation}\label{gu:eq:shalika-embedding}
 m_B(a)=m(a,\lambda_B(a)),\qquad
 S_B^{\GU}=m_B(K)N(X)\simeq K\ltimes N(X).
\end{equation}
On $\operatorname{Herm}_2(E)$ the map $u\mapsto aua^\dagger$ has absolute
determinant $|\det a|_E^2$. For diagonal $a$, this follows from its two
norm-scaled diagonal coordinates and its off-diagonal $E$-coordinate;
elementary unipotents have determinant one. Consequently
\begin{equation}\label{gu:jac:siegel-modulus}
 \delta_{P^{\GU}}(m(a,t))=|\det a|_E^2|t|_F^{-4}.
\end{equation}
Taking determinants in \eqref{gu:jac:character-stabilizer} gives
\begin{equation}\label{gu:jac:plane-determinant}
 N_{E/F}(\det a)=t^2,\qquad |\det a|_E=|t|_F^2.
\end{equation}
Hence $\delta_{P^{\GU}}|_K=1$, the group $K\ltimes N(X)$ is unimodular,
and normalized and unnormalized Jacquet actions agree on $K$.

For a smooth representation $\sigma$ of $K$, we use $\sigma\boxtimes\psi_B$ with the embedding $m_B$. Since
\[
 m_B(zI_2)=m(zI_2,N_{E/F}z)=zI_4,
\]
its central character on $E^\times$ is $\omega_\sigma$. In particular,
\begin{equation}\label{gu:eq:central-necessary}
 \Hom_{S_B^{\GU}}(\Pi,\sigma\boxtimes\psi_B)\ne0
 \quad\Longrightarrow\quad\omega_\Pi=\omega_\sigma.
\end{equation}

\subsection{The Schr\"odinger model}
On the underlying $F$-space of $V\otimes_EW$ put
\begin{equation}\label{gu:eq:tensor-symplectic}
 [v\otimes w,v'\otimes w']
 =-\tfrac12\Tr_{E/F}\bigl(h_V(v,v')h_W(w,w')\bigr).
\end{equation}
This is a nondegenerate alternating form, with polarization $(V\otimes X)\oplus(V\otimes Y)$. The second Lagrangian is identified with
\begin{equation}\label{gu:eq:schrodinger-Z}
 Z=Y\otimes_EV\simeq\Hom_E(X^c,V),
 \qquad
 (y\otimes v)(x\otimes1)=\overline{h_W(x,y)}\,v.
\end{equation}
It has dimension $16$ over $F$. Define the form-valued moment map by
\begin{equation}\label{gu:eq:moment-map}
 Q(T)=T^*h_V,\qquad G_T=T^\dagger J_4T.
\end{equation}
Thus $G_T$ is the matrix of the Hermitian form $Q(T)$ on $X^c$. The corresponding form matrix on $X$ is $\overline{G_T}$. The notation $T^*h_V$ means pullback of a form. If one instead forms an adjoint using $B^c$, the equation $Q(T)=cB^c$ becomes $T^*T=cI_{X^c}$.

For a locally compact totally disconnected space $Z'$, let $\mathcal S(Z')=C_c^\infty(Z',\mathbb C)$. The ordinary oscillator representation acts on $\mathcal S(Z)$. We use Heisenberg multiplication
\[
 (z,t)(z',t')=(z+z',t+t'+\tfrac12[z,z'])
\]
and Schr\"odinger operators
\begin{equation}\label{gu:eq:heisenberg-action}
 [\rho_\psi(x+y,t)f](T)
 =\psi\bigl(t-[x,T]-\tfrac12[x,y]\bigr)f(T+y)
\end{equation}
for $x\in V\otimes X$ and $y,T\in Z$. Weil operators satisfy left covariance
\[
 \omega(s)\rho_\psi(z,t)\omega(s)^{-1}=\rho_\psi(sz,t).
\]
The choices in~\eqref{gu:eq:tensor-symplectic} and~\eqref{gu:eq:heisenberg-action} give~\eqref{gu:eq:psiB}. Initially both ordinary splitting characters are one, which is permissible because both dimensions are even. Splitting changes will be introduced after the normalized calculation.

\subsection{Induction, quotients, and the unitary multiplicity-one theorem}
For a closed subgroup $D\subset L$ and a smooth representation $\rho$ of $D$, we realize the unnormalized compact induction $\cInd_D^L\rho$ as right-smooth functions satisfying
\[
 f(d\ell)=\rho(d)f(\ell),\qquad
 [\pi(\ell_0)f](\ell)=f(\ell\ell_0),
\]
with support compact modulo $D$. Ordinary smooth induction $\Ind_D^L\rho$ uses the same covariance without a compact-support condition. For a smooth representation $M$ the twisted Jacquet module is the unnormalized quotient
\[
 M_{N,\psi_B}
 =M/\langle nv-\psi_B(n)v:n\in N,\ v\in M\rangle.
\]

For $V=V_1\perp V_2$ with hyperbolic planes $V_1,V_2$, we retain
$H_{\mathrm U}$ for the equal-multiplier subgroup
\eqref{eq:main-unitary-subgroup}, to which
Theorem~\ref{thm:main-unitary} applies.

\subsection{The matching subgroup and its normalized extension}
\label{gu:jac:oscillator-section}
The direct tensor action has multiplier $\lambda(g)\lambda(h)$. We use the
equal-multiplier convention
\[
 R=\{(g,h)\in G_\theta\times G'_\theta:\lambda(g)=\lambda(h)\},\qquad
 \alpha_W(h)=\lambda(h)^{-1}h,
\]
and the symplectic action
\begin{equation}\label{gu:jac:tensor-action}
 (g,h)\longmapsto g\otimes\alpha_W(h).
\end{equation}
The involution $\alpha_W$ identifies this with the inverse-multiplier
convention and is the identity on $(G'_\theta)^1$. The tensor kernel is
\begin{equation}\label{gu:jac:tensor-kernel}
 K_E=\{(zI_V,\bar zI_W):z\in E^\times\},
\end{equation}
since $\alpha_W(\bar zI_W)=z^{-1}I_W$.
We use the similitude oscillator formalism of~\cite[\S\S2--4]{ZhangSimilitude}
and~\cite[\S\S3, 6.1]{BGS}. The following formulas fix its normalization.

For the ordinary Weil representation $\omega$ with splitting characters
$(1,1)$, the formulas needed below are
\begin{align}
 [\omega(g,1)f](T)&=f(g^{-1}T), &&g\in U(V),
 \label{gu:jac:ordinary-g}\\
 [\omega(1,m(a,1))f](T)&=|\det a|_E^2f(T\bar a),
 &&a\in\operatorname{GL}_2(E),\label{gu:jac:ordinary-levi}\\
 [\omega(1,n(u))f](T)&=
 \psi\left(-\frac12\operatorname{tr}_E
                     (u\overline{G_T})\right)f(T),
 &&u^\dagger=u.\label{gu:jac:ordinary-n}
\end{align}
These are the ordinary unitary formulas of \cite[\S3.1]{ChenGan} in our matrix
and trace conventions.

Their phase, including its sign and conjugation, follows from
\eqref{gu:eq:heisenberg-action}; the calculation and the Fourier operators
are recorded in Appendix~\ref{gu:app:normalization}.

For $s\in F^\times$, put
\begin{equation}\label{gu:jac:sections}
 r_s=\operatorname{diag}(s,s,1,1)\in G_\theta,
 \qquad d_s=\operatorname{diag}(I_2,sI_2)\in G'_\theta,
 \qquad e_s=(r_s,d_s).
\end{equation}
Both $r_s$ and $d_s$ have multiplier $s$, and these are multiplicative
sections.  If $(g,h)\in R$ has common multiplier $s$, its unique
decomposition is
\begin{equation}\label{gu:jac:semidirect-decomposition}
 (g,h)=(gr_s^{-1},hd_s^{-1})e_s,
 \qquad (gr_s^{-1},hd_s^{-1})\in U(V)\times U(W).
\end{equation}
Define
\begin{equation}\label{gu:jac:dilation}
 [D_sf](T)=|s|_F^{-4}f(r_s^{-1}T).
\end{equation}
The map $T\mapsto r_sT$ has absolute $F$-determinant $|s|_F^8$, since
$|\det_Er_s|_E=|s|_F^4$ and $T$ has two columns.  Thus the coefficient
in \eqref{gu:jac:dilation} is its inverse positive square root.  In
particular, $D_sD_t=D_{st}$, and $D_s$ is unitary for the Schr\"odinger
$L^2$-measure.

\begin{proposition}\label{gu:jac:extension}
The formula
\begin{equation}\label{gu:jac:extension-definition}
 \omega_{R,0}(g,h)
 =\omega(gr_s^{-1},hd_s^{-1})D_s,
 \qquad s=\lambda(g)=\lambda(h),
\end{equation}
defines a smooth oscillator representation of $R$ implementing
\eqref{gu:jac:tensor-action} and extending $\omega$.  Its action on the
matching Siegel Levi is
\begin{equation}\label{gu:jac:normalized-levi}
 [\omega_{R,0}(g,m(a,s))f](T)
 =|s|_F^{-4}|\det a|_E^2f(g^{-1}T\iota(a)),
 \qquad s=\lambda(g).
\end{equation}
The subgroup $K_E$ acts trivially.
\end{proposition}

\begin{proof}
The semidirect decomposition reduces the representation law to
\begin{equation}\label{gu:jac:covariance}
 D_s\omega(g_0,h_0)D_s^{-1}
 =\omega(r_sg_0r_s^{-1},d_sh_0d_s^{-1}).
\end{equation}
The generator calculation is given in Appendix~\ref{gu:app:normalization}.
Substituting $hd_s^{-1}=m(a,1)$ in \eqref{gu:jac:extension-definition}
gives \eqref{gu:jac:normalized-levi}. Its coefficient is the positive square
root of the absolute determinant $|s|_F^{-8}|\det a|_E^4$ of
$T\mapsto g^{-1}T\bar a$. For $(zI_V,\bar zI_W)\in K_E$, the argument
is $T$ and the scalar is
$|N_{E/F}z|_F^{-4}|\det(\bar zI_2)|_E^2=1$.
\end{proof}

\subsection{The product group and its multiplier coordinate}
The quotient map and a multiplicative section are
\[
 q(g,h)=\lambda(h)/\lambda(g),\qquad q_c=(1,d_c).
\]
The decomposition $G_\theta\times G'_\theta=R\{q_c:c\in F^\times\}$ identifies
$R\backslash(G_\theta\times G'_\theta)$ homeomorphically with $F^\times$.
Both groups are unimodular: $R$ is reductive, with kernel
$G_\theta^1\times (G'_\theta)^1$ under its common multiplier. Set
\begin{equation}\label{gu:jac:full-product}
 \Omega_0^{\GU}=\cInd_R^{G_\theta\times G'_\theta}\omega_{R,0}.
\end{equation}
\begin{proposition}\label{gu:jac:full-model}
Evaluation along $q_c$ gives a realization
\[
 \Omega^{\GU}_0=\mathcal S(F^\times\times Z),
 \qquad \phi(c,T)=F(q_c)(T).
\]
For $(g,h)\in G_\theta\times G'_\theta$, with $s=\lambda(g)$ and $t=\lambda(h)$, its action is
\begin{equation}\label{gu:jac:induced-action}
 [\Omega^{\GU}_0(g,h)\phi](c,\cdot)
 =\omega_{R,0}\bigl(g,d_chd_{ct/s}^{-1}\bigr)
                   \phi(ct/s,\cdot).
\end{equation}
In particular,
\begin{align}
 [\Omega^{\GU}_0(1,n(u))\phi](c,T)
 &=\psi\left(-\frac1{2c}\operatorname{tr}_E
                            (u\overline{G_T})\right)\phi(c,T),
 \label{gu:jac:induced-n}\\
 [\Omega^{\GU}_0(g,m(a,t))\phi](c,T)
 &=|s|_F^{-4}|\det a|_E^2
                         \phi(ct/s,g^{-1}T\bar a).
 \label{gu:jac:induced-levi}
\end{align}
\end{proposition}

\begin{proof}
Evaluation and the inverse $F(rq_c)=\omega_{R,0}(r)\phi(c,\cdot)$
identify the spaces: compactly supported section values have a finite
compact open constancy partition, and smoothness supplies a common
small right stabilizer. These maps also preserve the Schwartz
inductive-limit topologies. The identity
\[
 q_c(g,h)=(g,d_chd_{ct/s}^{-1})q_{ct/s}
\]
gives \eqref{gu:jac:induced-action}. Now
$d_cn(u)d_c^{-1}=n(c^{-1}u)$ and
$d_cm(a,t)d_{ct/s}^{-1}=m(a,s)$ give the two asserted formulas.
\end{proof}
The central subgroup $K_E$ acts trivially on $\Omega_0^{\GU}$ as well.



\subsection{The scaled Gram orbit}\label{gu:jac:jacquet-section}
Fix an isometric embedding $j:(X^c,B^c)\hookrightarrow V$, and put
$L=G_\theta\times K$. Define
\begin{equation}\label{gu:jac:scaled-locus}
 \widetilde{\mathcal O}_B
 =\{(c,T)\in F^\times\times Z:Q(T)=cB^c\}.
\end{equation}
The following is the unitary similitude analogue of
\cite[Proposition~3.4]{ChenGan}.

\begin{proposition}[Similitude Jacquet module]\label{gu:jac:jacquet}
Restriction, followed by multiplication by $|c|_F^4$, gives isomorphisms
\begin{equation}\label{gu:jac:J}
 (\Omega_0^{\GU})_{N(X),\psi_B}
 \simeq\mathcal S(\widetilde{\mathcal O}_B)
 \simeq\cInd_{H_j}^{G_\theta\times K}\mathbf1,
\end{equation}
where the middle term has geometric pullback action for
\begin{equation}\label{gu:jac:orbit-action}
 (g,a)\cdot(c,T)
 =\left(c\frac{\lambda(g)}{\lambda_B(a)},\ gT\iota(a)^{-1}\right),
\end{equation}
and
\begin{equation}\label{gu:jac:stabilizer}
 H_j=\{(g,a):\lambda(g)=\lambda_B(a),\ gj=j\iota(a)\}.
\end{equation}
These are homeomorphisms for the Schwartz inductive-limit topologies.
In these coordinates, the composite is
\begin{equation}\label{gu:jac:intertwiner}
 \mathcal I_j[\phi](g,a)
 =\left|\frac{\lambda_B(a)}{\lambda(g)}\right|_F^4
 \phi\left(\frac{\lambda_B(a)}{\lambda(g)},g^{-1}j\iota(a)\right).
\end{equation}
\end{proposition}

\begin{proof}
\emph{Twisted support.}
By \eqref{gu:jac:induced-n} and the nondegenerate trace pairing from
the conventions above, the unipotent character equals $\psi_B$ exactly
when $c^{-1}\overline{G_T}=\mathsf B$, or $Q(T)=cB^c$.
The locus is closed, and every $T$ on it is injective with nondegenerate
image. Lemma~\ref{gu:jac:restriction} gives the surjectivity of
restriction, identifies its kernel with the algebraic Jacquet relation
space, and establishes the quotient topology.

For $s=\lambda(g)$ and $t=\lambda_B(a)$, equations
\eqref{gu:jac:induced-levi} and \eqref{gu:jac:plane-determinant} give the
action before changing realization:
\begin{equation}\label{gu:jac:quotient-action}
 [\Omega_0^{\GU}(g,a)\phi](c,T)
 =|t/s|_F^4\phi(ct/s,g^{-1}T\iota(a)).
\end{equation}
The Gram matrix at the new argument is
$s^{-1}a^tG_T\bar a=(ct/s)\overline{\mathsf B}$, by
\eqref{gu:jac:character-stabilizer}. Set
\begin{equation}\label{gu:jac:gauge}
 \Phi(c,T)=|c|_F^4\phi(c,T).
\end{equation}
Since $|c|_F^4|t/s|_F^4=|ct/s|_F^4$, this gives the pullback action $\pi_0$
for \eqref{gu:jac:orbit-action}; multiplication preserves support and the
Schwartz topology.

\emph{Rational orbit.}
For $(c,T)$ in the locus, $r_c^{-1}T$ is an isometric embedding.
Witt extension~\cite[Chapter~7, \S9]{Scharlau85} gives $u\in U(V)$ with $uj=r_c^{-1}T$; hence
$(r_cu,1)\cdot(1,j)=(c,T)$.
The orbit map is a submersion and induces a homeomorphism
$L/H_j\simeq\widetilde{\mathcal O}_B$, by
Lemma~\ref{gu:jac:orbit-topology} in the appendix.

\emph{Stabilizer and block transport.}
At $o=(1,j)$, equation \eqref{gu:jac:orbit-action} gives
\eqref{gu:jac:stabilizer}. Put
\begin{equation}\label{gu:jac:planes}
 V_2=j(X^c),\qquad V_1=V_2^\perp,\qquad
 t_j(a)=j\iota(a)j^{-1}.
\end{equation}
Both planes are hyperbolic, by Witt extension and orthogonal
complementation in the split four-space. The map
$t_j:K\xrightarrow{\sim}\GU(V_2)$ preserves multipliers.
A similitude preserving $V_2$ also preserves $V_1$, and therefore
\begin{equation}\label{gu:jac:block-stabilizer}
 H_j=\{(u\oplus h,t_j^{-1}(h)):
       \lambda(u)=\lambda(h),\ u\in\GU(V_1),\ h\in\GU(V_2)\}.
\end{equation}
Transport by $1\times t_j$ identifies it with the subgroup
\eqref{eq:main-unitary-subgroup}.

\emph{Fiber character and modulus.}
For a stabilizer element of common multiplier $s$, the matching Levi
operator at $c=1$ has fixed argument and scalar
\begin{equation}\label{gu:jac:trivial-fiber}
 |s|_F^{-4}|\det a|_E^2=1.
\end{equation}
The splitting characters and the multiplier normalization are those of
Proposition~\ref{gu:jac:extension}, so the fiber character is trivial.
The equal-multiplier product $H_j$ is reductive, with kernel
$U(V_1)\times U(V_2)$ under its multiplier. Thus $H_j$ and $L$ are
unimodular and
\begin{equation}\label{gu:jac:relative-modulus}
 \Delta_L|_{H_j}\Delta_{H_j}^{-1}=1.
\end{equation}
The elements
$(zI_V,\bar zI_X)$ of $H_j$ transport to the common diagonal
$E^\times$-center of the block subgroup.

\emph{Compact induction.}
Compose restriction and \eqref{gu:jac:gauge} with
$[\mathcal A\Phi](k)=\Phi(k^{-1}o)$. The quotient homeomorphism
identifies compact support with compact support modulo $H_j$, and gives
the inverse. For $h\in H_j$ and $k_0\in L$,
\[
 \mathcal A\Phi(hk)=\mathcal A\Phi(k),\qquad
 \mathcal A(\pi_0(k_0)\Phi)(k)=\mathcal A\Phi(kk_0).
\]
The composite is \eqref{gu:jac:intertwiner}.
\end{proof}

Changing $j$ to $vj$, with $v\in U(V)$, conjugates $H_j$ by $(v,1)$
and $t_j$ by $v$, giving the natural isomorphism of the induced models.
On restriction to $U(V)\times U(W)$, the coordinate $c$ remains;
its $c=1$ fiber recovers the ordinary Weil-model calculation of
\cite[Proposition~3.4]{ChenGan}.


\subsection{Splitting characters and the Hom identity}\label{gu:sec:hom}
For $\tau\in\Irr(G_\theta)$ define the big theta lift by
\begin{equation}\label{gu:eq:big-theta-zero}
 \Theta_0^{\GU}(\tau)=(\Omega_0^{\GU}\otimes\tau^\vee)_{G_\theta}.
\end{equation}
Lemma~\ref{gu:app:comparison} identifies this coinvariant construction with the matching-subgroup construction, with the forward
integral \eqref{gu:eq:forward-comparison}.

We use the family of splitting data $D=(\chi,1,1)$, where $\chi$ is a
unitary character of $E^\times/F^\times$. Let
\[
 \mathfrak i:E^\times/F^\times\xrightarrow{\sim}E^1,
 \quad z\mapsto z/\bar z,\qquad \mu_\eta=\eta\circ\mathfrak i^{-1}.
\]
For either four-dimensional similitude group set
\begin{equation}\label{gu:eq:ambient-character}
 d_4(g)=\frac{\det_Eg}{\lambda(g)^2}\in E^1,\qquad
 \mathcal A_\eta(g)=\mu_\eta(d_4(g)).
\end{equation}
The determinant relation $N_{E/F}(\det_Eg)=\lambda(g)^4$ makes this a
character on either similitude group. Its isometry restriction gives the
splitting change in~\cite[equations~(3.1)--(3.3)]{ChenGan}: on $U(V)$ it is
$\mu_\eta\circ\det_E$, and on $m(a,1)$ it is $\eta(\det_Ea)$.
It is trivial on the upper unipotent radical and on the skew-Hermitian
partial Weyl generators, whose determinants are one. Put
\begin{equation}\label{gu:eq:chi-oscillator}
 \omega_{R,\chi}=\omega_{R,0}(1\boxtimes\mathcal A_\chi)|_R,
 \qquad \Omega_\chi^{\GU}=\cInd_R^{G_\theta\times G'_\theta}\omega_{R,\chi},
 \qquad \Theta_{\chi,1,1}^{\GU}(\tau)
       =(\Omega_\chi^{\GU}\otimes\tau^\vee)_{G_\theta}.
\end{equation}
Here $\chi=1$ recovers the subscript $0$ notation. Multiplying inducing
functions by the ambient character gives
\begin{equation}\label{gu:eq:chi-twist}
 \Omega_\chi^{\GU}\simeq
 \Omega_0^{\GU}\otimes(1\boxtimes\mathcal A_\chi),\qquad
 \Theta_{\chi,1,1}^{\GU}(\tau)
 \simeq\mathcal A_\chi\otimes\Theta_0^{\GU}(\tau).
\end{equation}
The general data $(\chi_V,\chi_W,\kappa)$ are treated in
Appendix~\ref{gu:app:splitting}.

\begin{theorem}[Shalika period identity for similitude theta lifts]\label{gu:thm:hom}
For irreducible admissible $\tau$ of $G_\theta$ and $\sigma$ of $K$, and any
unitary character $\chi$ of $E^\times/F^\times$, there is an isomorphism
\begin{equation}\label{gu:eq:hom-coverage}
 \begin{aligned}
 &\Hom_{S_B^{\GU}}(\Theta_{\chi,1,1}^{\GU}(\tau),
                         \sigma\boxtimes\psi_B)\\
 &\qquad\simeq
 \Hom_{H_{\mathrm U}}\left(\tau^\vee\boxtimes
  \bigl((\sigma^j)^\vee\otimes(\chi^{-1}\circ\det_E)\bigr),
  \mathbf1\right),\qquad \sigma^j=\sigma\circ t_j^{-1}.
 \end{aligned}
\end{equation}
Its scalar normalization is fixed by a nonzero invariant measure on
$H_j\backslash(G_\theta\times K)$. In particular, the Hom space has dimension
at most one.
\end{theorem}
\begin{proof}
First take $\chi=1$. Tensor evaluation (Lemma~\ref{gu:app:tensor}),
Proposition~\ref{gu:jac:jacquet}, and closed-subgroup duality
(Lemma~\ref{gu:app:duality}) give the canonical chain
\begin{equation}\label{gu:eq:hom-zero}
 \begin{aligned}
 \Hom_{S_B^{\GU}}(\Theta_0^{\GU}(\tau),\sigma\boxtimes\psi_B)
 &\simeq\Hom_{G_\theta\times K}((\Omega_0^{\GU})_{N,\psi_B},
                                      \tau\boxtimes\sigma)\\
 &\simeq\Hom_{G_\theta\times K}(\cInd_{H_j}^{G_\theta\times K}\mathbf1,
                                      \tau\boxtimes\sigma)\\
 &\simeq\Hom_{H_j}(\tau^\vee\boxtimes\sigma^\vee,\mathbf1).
 \end{aligned}
\end{equation}
In the first identification, tensor evaluation is followed by factoring
through the $(N,\psi_B)$-quotient.
The last step applies because $H_j$ and $G_\theta\times K$ are unimodular and
$\tau\boxtimes\sigma$ is admissible by~\cite[Proposition~2.16]{BZ76}.
Transporting the coefficient by $t_j$ gives the corresponding
spherical-pair Hom space.

For general $\chi$, on the Shalika subgroup
\[
 d_4(m_B(a))=\mathfrak i(\det_Ea),\qquad
 \mathcal A_\chi(m_B(a))=\chi(\det_Ea),\qquad
 \mathcal A_\chi|_{N(X)}=1.
\]
After absorbing this target character through \eqref{gu:eq:chi-twist},
the normalized identity gives the following three changes of coefficient:
\begin{equation}\label{gu:eq:coefficient-three-steps}
 \begin{aligned}
 \sigma_0&=\sigma\otimes(\chi^{-1}\circ\det_X),\\
 \sigma_0^j&=\sigma^j\otimes(\chi\circ\det_{V_2}),\\
 (\sigma_0^j)^\vee&=(\sigma^j)^\vee
                        \otimes(\chi^{-1}\circ\det_{V_2}).
 \end{aligned}
\end{equation}
Indeed, $t_j$ conjugates determinants by \eqref{gu:eq:iota-det}, and
$\chi(\bar z)=\chi(z)^{-1}$. This proves \eqref{gu:eq:hom-coverage}.
The representation on its right is irreducible admissible on the
specified split product, so Theorem~\ref{thm:main-unitary} gives the bound.
\end{proof}
